\underline{}% ----------------------------------------------------------
% Resumo
% ----------------------------------------------------------
\chapter*{Resumo}\label{resumo}
% ----------------------------------------------------------

%Ranqueamento em Classificadores Associativos através da Agregação de Medidas Objetivas

A mineração de dados é o processo de descobrir padrões e conhecimentos a partir de grandes quantidades de dados objetivando analisar padrões frequentes e correlações contidas nos dados, com especial interesse em padrões inesperados, envolvendo assim um conjunto de métodos e técnicas. Apesar dessas técnicas já estarem bem difundidas em diversos setores da sociedade, percebe-se grande potencial de utilização a ser desenvolvido no setor público, em especial na área de orçamentação pública.  
Neste contexto, este trabalho se propõe a avaliar o uso de técnicas de inteligência computacional no campo do orçamento público de forma geral e, especificamente, na tarefa específica de previsão das despesas orçamentárias.

\textbf{Palavras-chaves:} mineração de dados; aprendizado de máquina; orçamento público.

%\par


%\chapter*{Abstract}\label{Abstract}
% ----------------------------------------------------------


%\textbf{Keywords:} associative classifiers; school dropout; imbalanced data; data mining.

\par
